% Paper 4, §7 — Sparse autoencoders: the unsupervised feature framing
% Thesis: SAEs commit to features as sparse linear directions in residual
% space. TopK (Gao 2024) and JumpReLU (Rajamanoharan 2024) are the
% production-grade variants; Gemma Scope is the open-suite scale
% demonstration; transcoders bridge the framing into circuit analysis.

\section{Sparse autoencoders}
\label{sec:sparse-autoencoders}

A sparse autoencoder (SAE) makes a specific commitment about what a
feature \emph{is}: a direction $\mathbf{d}_i \in \mathbb{R}^n$ in
activation space such that, on any given input, the activation
$\mathbf{x}$ is well-approximated by a sparse linear combination of a
small subset of such directions drawn from an overcomplete dictionary
$\{\mathbf{d}_i\}_{i=1}^M$ with $M \gg n$. The commitment is strong —
features are linear, additive, and identifiable from activations alone
without supervision — and the methodological payoff is that activations
acquire an addressable, named basis on which downstream interpretability
work (probing, patching, circuit tracing) can land. This section covers
the two production-grade architectures (Top-K \citep{gao2024-topk-saes},
JumpReLU \citep{rajamanoharan2024-jumprelu}), the open scale
demonstration (Gemma Scope \citep{gemma-scope-2024}), and the structural
cousin (transcoders \citep{dunefsky2024-transcoders}) that carries the
framing into the circuit-tracing pipeline of \S\ref{sec:circuits}. The
canonicality question raised by these commitments is taken up in
\S\ref{sec:sae-canonical}.

\subsection{The framing: dictionary learning on activations}
\label{sec:sae-framing}

Given an activation $\mathbf{x} \in \mathbb{R}^n$ — typically the
residual stream at some layer $\ell$, or an MLP/attention-output vector
— an SAE is a pair of an encoder and a decoder
\citep[\S 2, Eq.~1--2]{rajamanoharan2024-jumprelu}\footnote{KB excerpt:
\texttt{kb/excerpts/rajamanoharan2024-jumprelu.md\#sec-2}.}:
\begin{align}
  \mathbf{f}(\mathbf{x}) &:= \sigma\!\left(\mathbf{W}_{\text{enc}}\,\mathbf{x} + \mathbf{b}_{\text{enc}}\right),
  \label{eq:sae-encoder}\\
  \hat{\mathbf{x}}(\mathbf{f}) &:= \mathbf{W}_{\text{dec}}\,\mathbf{f} + \mathbf{b}_{\text{dec}},
  \label{eq:sae-decoder}
\end{align}
where $\mathbf{W}_{\text{enc}} \in \mathbb{R}^{M \times n}$,
$\mathbf{W}_{\text{dec}} \in \mathbb{R}^{n \times M}$, and the encoder
activation $\sigma$ is the architectural choice that defines the SAE
family. The columns of $\mathbf{W}_{\text{dec}}$, written
$\mathbf{d}_i$, are the \emph{decoder feature directions}; the
expansion factor $M / n$ is typically $8$--$1000$. The training
objective always combines reconstruction and sparsity:
\begin{equation}
  \mathcal{L}(\mathbf{x})
  \;=\;
  \underbrace{\|\mathbf{x} - \hat{\mathbf{x}}(\mathbf{f}(\mathbf{x}))\|_2^2}_{\mathcal{L}_{\text{reconstruct}}}
  \;+\;
  \mathcal{L}_{\text{sparsity}}.
  \label{eq:sae-loss}
\end{equation}
The sparsity term and the activation $\sigma$ jointly determine the
position the SAE occupies on the $L_0$-vs-reconstruction Pareto
frontier, and that frontier is, empirically, the only training axis
that matters: every modern SAE family is an attempt to push the
Pareto curve down-and-left at fixed compute
\citep[\S 1]{rajamanoharan2024-jumprelu}.

The original recipe of \citet{templeton2024-scaling-monosemanticity}
following \citet{gao2024-topk-saes}'s baseline uses ReLU encoder with
an $\ell_1$ sparsity penalty $\lambda \|\mathbf{f}\|_1$
\citep[\S 2.2, Eq.~1]{gao2024-topk-saes}\footnote{KB excerpt:
\texttt{kb/excerpts/gao2024-topk-saes.md\#sec-2-2-baseline}.}. Two
known issues motivate the TopK and JumpReLU successors. First,
\emph{shrinkage}: $\ell_1$ biases all positive activations toward zero,
so reconstructions systematically underestimate true feature
magnitudes \citep[\S 2.3]{gao2024-topk-saes}. Second,
\emph{reparameterization non-invariance}: scaling encoder rows up and
decoder columns down preserves $\mathbf{W}_{\text{dec}}\mathbf{f}$ but
changes $\|\mathbf{f}\|_1$, so the $\ell_1$ penalty is degenerate
without a unit-norm constraint on $\mathbf{d}_i$
\citep[\S 3]{rajamanoharan2024-jumprelu}.

\subsection{Top-K SAEs (Gao et al.\ 2024)}
\label{sec:sae-topk}

\citet{gao2024-topk-saes} replace the ReLU+$\ell_1$ recipe with a
$k$-sparse activation \citep[after][]{gao2024-topk-saes}: keep the
top-$K$ pre-activations, zero the rest. The encoder becomes
\citep[\S 2.3, Eq.~2]{gao2024-topk-saes}\footnote{KB excerpt:
\texttt{kb/excerpts/gao2024-topk-saes.md\#sec-2-3}.}
\begin{equation}
  \mathbf{f}(\mathbf{x}) \;=\; \mathrm{TopK}\!\left(\mathbf{W}_{\text{enc}}(\mathbf{x} - \mathbf{b}_{\text{pre}})\right),
  \label{eq:topk}
\end{equation}
and because the sparsity is enforced exactly --- $L_0 = K$ at every
token by construction --- the loss reduces to
$\mathcal{L} = \|\mathbf{x} - \hat{\mathbf{x}}\|_2^2$ with the $\ell_1$
penalty removed entirely. The four claimed benefits are direct
consequences of this design choice
\citep[\S 2.3]{gao2024-topk-saes}: (i) no shrinkage bias from $\ell_1$;
(ii) $L_0$ is set directly rather than tuned via $\lambda$, simplifying
sweeps and model comparison; (iii) the empirical sparsity-reconstruction
frontier improves over ReLU+$\ell_1$, with the gap widening at scale
\citep[Fig.~2b]{gao2024-topk-saes}; (iv) clamping small activations to
zero increases the monosemanticity of random activating examples
\citep[\S 4.3]{gao2024-topk-saes}.

The cost is the \emph{dead-latent problem}: at large $M$, latents that
fail to fire early in training never recover. \citet{gao2024-topk-saes}
report that without mitigation, ``up to 90\% dead
latents''~\citep[\S 2.4]{gao2024-topk-saes} appear at the largest
scales\footnote{KB excerpt:
\texttt{kb/excerpts/gao2024-topk-saes.md\#sec-2-4}.}. Two ingredients
fix it: encoder rows are initialized as the transpose of decoder
columns; and an auxiliary ``AuxK'' loss models the residual
reconstruction error using the top-$k_{\text{aux}}$ \emph{dead}
latents, forcing them to reactivate. With these, even a
$16$M-latent SAE on GPT-4 residual-stream activations trained for $40$B
tokens keeps only $7\%$ of latents
dead~\citep[\S 1, \S 2.4]{gao2024-topk-saes}\footnote{KB excerpt:
\texttt{kb/excerpts/gao2024-topk-saes.md\#sec-1-headline}.}. Optimal
compute follows a clean power law $L(C) = 0.09 + 0.056\,C^{-0.084}$ on
the MSE-vs-compute frontier \citep[Fig.~1, \S 3]{gao2024-topk-saes}, so
the top-K family scales predictably to GPT-4-class models.

\subsection{JumpReLU SAEs (Rajamanoharan et al.\ 2024)}
\label{sec:sae-jumprelu}

\citet{rajamanoharan2024-jumprelu} take a different route to the same
goal: replace ReLU with a learned-threshold step,
\begin{equation}
  \mathrm{JumpReLU}_{\boldsymbol{\theta}}(z) \;:=\; z\,H(z - \theta),
  \label{eq:jumprelu}
\end{equation}
where $H$ is the Heaviside step function and
$\boldsymbol{\theta} \in \mathbb{R}^M_+$ is a \emph{per-feature learned
threshold} below which encoder pre-activations are set to zero
\citep[\S 2, Eq.~4]{rajamanoharan2024-jumprelu}\footnote{KB excerpt:
\texttt{kb/excerpts/rajamanoharan2024-jumprelu.md\#sec-2-activations}.}.
The full SAE encoder is
$\mathbf{f}(\mathbf{x}) := \mathrm{JumpReLU}_{\boldsymbol{\theta}}(\mathbf{W}_{\text{enc}}\mathbf{x} + \mathbf{b}_{\text{enc}})$
\citep[\S 3, Eq.~8]{rajamanoharan2024-jumprelu}, and sparsity is
trained directly via an $L_0$ penalty rather than its $\ell_1$ proxy:
\begin{equation}
  \mathcal{L}(\mathbf{x})
  \;=\;
  \|\mathbf{x} - \hat{\mathbf{x}}(\mathbf{f}(\mathbf{x}))\|_2^2
  \;+\;
  \lambda\,\|\mathbf{f}(\mathbf{x})\|_0.
  \label{eq:jumprelu-loss}
\end{equation}
Equation~\eqref{eq:jumprelu-loss} is non-differentiable both in
$\boldsymbol{\theta}$ (because of $H$) and in $\|\cdot\|_0$, which
\citet{rajamanoharan2024-jumprelu} resolve with
\emph{straight-through estimators}: the backward pass replaces the
discontinuous derivative with a kernel-density-estimated surrogate of
the derivative of the expected loss, separately for $\boldsymbol{\theta}$
and for the $L_0$ term \citep[\S 3]{rajamanoharan2024-jumprelu}. The
key design point is that this avoids the shrinkage bias of $\ell_1$
entirely while keeping the encoder elementwise --- no partial sort, no
batchwise communication --- so a JumpReLU SAE trains as cheaply as a
plain ReLU SAE.

The empirical headline: at any fixed sparsity level on Gemma 2 9B
residual-stream, MLP-output, and attention-output activations, JumpReLU
SAEs ``provide reconstructions that are at least as good as, and often
slightly better than, TopK SAEs'' and consistently dominate Gated SAEs
\citep[\S 1]{rajamanoharan2024-jumprelu}\footnote{KB excerpt:
\texttt{kb/excerpts/rajamanoharan2024-jumprelu.md\#sec-1-pareto}.}. The
caveat is interpretability, not reconstruction: both TopK and JumpReLU
produce a small fraction of \emph{high-frequency} latents (firing on
$>10\%$ of tokens, fewer than $0.06\%$ of a $131$K-width SAE) which
manual ratings find systematically less interpretable
\citep[\S 1]{rajamanoharan2024-jumprelu}\footnote{KB excerpt:
\texttt{kb/excerpts/rajamanoharan2024-jumprelu.md\#sec-1-caveat}.}.
Pareto-winning on reconstruction-at-fixed-sparsity does not imply
Pareto-winning on per-feature interpretability.

\begin{intuition}
A sparse autoencoder is a search for the directions in activation space
such that most points use only $K$ of them. The encoder asks ``which
$K$ dictionary atoms best explain this activation?''; the decoder
reassembles the activation as a $K$-sparse linear combination
$\hat{\mathbf{x}} = \sum_{i \in S(\mathbf{x})} f_i\,\mathbf{d}_i$ with
$|S(\mathbf{x})| = K$. The choice of activation $\sigma$ in
Eq.~\eqref{eq:sae-encoder} is the choice of \emph{how} to enforce that
$K$: TopK does it by argmax (exactly $K$ active per token); JumpReLU
does it by a learned per-feature threshold (approximately $K$ on
average, learned). Both reduce to the same canonical form
$\hat{\mathbf{x}} = \sum_i f_i\,\mathbf{d}_i$ with $\mathbf{f}$
sparse --- the dictionary-learning math of
Eq.~\eqref{eq:sae-encoder}--\eqref{eq:sae-decoder}.
\end{intuition}

\subsection{Gemma Scope: the open scale demonstration}
\label{sec:sae-gemma-scope}

Until 2024, large-scale SAE work was largely Anthropic-internal
(\citealp{templeton2024-scaling-monosemanticity} on Claude 3 Sonnet)
or OpenAI-internal (\citealp{gao2024-topk-saes} on GPT-4). \emph{Gemma
Scope} \citep{gemma-scope-2024} is the open-release inflection: a
suite of JumpReLU SAEs trained on \emph{all} layers and sublayers ---
residual stream, MLP output, attention output --- of Gemma 2 2B and
9B base models, plus selected layers of Gemma 2 27B. The release
contains ``more than 400 sparse autoencoders in the main release, with
more than 30 million learned features in total''
\citep[\S 1]{gemma-scope-2024}\footnote{KB excerpt:
\texttt{kb/excerpts/gemma-scope-2024.md\#sec-1}.}, trained on
$4$--$16$B tokens each at widths $\{16.4\text{K}, 65.5\text{K},
131\text{K}, 1\text{M}\}$. The compute headline:
\begin{quote}
\itshape
We used over 20\% of the training compute of GPT-3 (Brown et al., 2020),
saved about 20 Pebibytes (PiB) of activations to disk, and produced
hundreds of billions of sparse autoencoder parameters in
total~\citep[\S 1]{gemma-scope-2024}.
\end{quote}
The justification for breadth is methodological:
\citet[\S 1]{gemma-scope-2024} argue that
``a comprehensive suite is essential for enabling more ambitious
applications of interpretability, such as circuit analysis,'' which
require feature dictionaries at every site through which causal
information flows.

A load-bearing implementation detail for the cross-method work in
\S\ref{sec:circuits}: \citet[\S 4.1]{gemma-scope-2024} report that
delta-loss is consistently higher for residual-stream SAEs than for
MLP- or attention-sublayer SAEs at matched FVU, because
``even small errors in reconstructing the residual stream can have a
significant impact on LM loss, whereas relatively large errors in
reconstructing a single MLP or attention sub-layer's output have a
limited impact on the LM loss'' \citep[\S 4.1]{gemma-scope-2024}.
Residual SAEs sit on the model's communication bus and any
reconstruction error compounds; sublayer SAEs sit on additive
contributions and absorb error locally. Circuit-analysis work that
re-inserts SAE reconstructions in place of the original activations
must account for this: residual-SAE reconstructions are the
high-leverage failure mode.

The release explicitly disclaims a converged evaluation framework:
``as of now there is no consensus on what constitutes a reliable metric
for the quality of a sparse autoencoder or its learned latents and
[\dots] this is an ongoing area of research and debate''
\citep[\S 4]{gemma-scope-2024}. Gemma Scope reports several metrics
(delta-loss, FVU, downstream-task evaluation) precisely because no
single one is canonical --- a frontier point we develop in
\S\ref{sec:sae-canonical}.

\subsection{Transcoders: from feature decomposition to circuit primitive}
\label{sec:sae-transcoders}

A \emph{transcoder} \citep{dunefsky2024-transcoders} is structurally an
SAE but trained on a different objective: rather than reconstructing an
activation, it learns to mimic an MLP sublayer's input-to-output
mapping with a sparse hidden representation in between. The
architecture
\citep[\S 3.1, Eq.~3--5]{dunefsky2024-transcoders}\footnote{KB excerpt:
\texttt{kb/excerpts/dunefsky2024-transcoders.md\#sec-3-1}.} is:
\begin{align}
  \mathbf{z}_{\text{TC}}(\mathbf{x})
  &= \mathrm{ReLU}\!\left(\mathbf{W}_{\text{enc}}\,\mathbf{x} + \mathbf{b}_{\text{enc}}\right),
  \label{eq:transcoder-encoder}\\
  \mathrm{TC}(\mathbf{x})
  &= \mathbf{W}_{\text{dec}}\,\mathbf{z}_{\text{TC}}(\mathbf{x}) + \mathbf{b}_{\text{dec}},
  \label{eq:transcoder-decoder}
\end{align}
trained against the MLP function the transcoder is meant to replace:
\begin{equation}
  \mathcal{L}_{\text{TC}}(\mathbf{x})
  \;=\;
  \underbrace{\|\mathrm{MLP}(\mathbf{x}) - \mathrm{TC}(\mathbf{x})\|_2^2}_{\text{faithfulness}}
  \;+\;
  \lambda_1\,\|\mathbf{z}_{\text{TC}}(\mathbf{x})\|_1.
  \label{eq:transcoder-loss}
\end{equation}
The structural difference from an SAE is one substitution in the
faithfulness term: an SAE-on-MLP-output minimizes
$\|\mathrm{MLP}(\mathbf{x}) - \hat{\mathbf{x}}\|_2^2$ where
$\hat{\mathbf{x}}$ depends on the \emph{post}-MLP activation, and so
the SAE's feature dictionary is a decomposition of that
\emph{specific} activation. A transcoder's $\mathrm{TC}(\mathbf{x})$
depends on the \emph{pre}-MLP input $\mathbf{x}$, and so its feature
dictionary is a decomposition of the MLP \emph{function}, not its
particular output \citep[\S 1]{dunefsky2024-transcoders}\footnote{KB
excerpt: \texttt{kb/excerpts/dunefsky2024-transcoders.md\#sec-1-comparison}.}.

The methodological consequence \citep[\S 1]{dunefsky2024-transcoders}:
the connection from a pre-MLP feature $f_i$ (associated with decoder
column $\mathbf{d}_i$) to a post-transcoder feature $f_j$ (associated
with encoder row $\mathbf{e}_j$) factorizes into an
\emph{input-dependent} term (the activations $f_i(\mathbf{x})$ and
$f_j(\mathbf{x})$ at the input of interest) and an \emph{input-invariant}
term (the bilinear weight product
$\mathbf{e}_j^\top\mathbf{d}_i$ that holds for any input). This is the
factorization that the cross-layer transcoder of
\citet{lindsey2025-circuit-tracing} generalizes to skip-layer feature
flow, and it is the architectural primitive on which the
attribution-graph construction of \S\ref{sec:circuit-attribution} is
built. As an analytical tool, transcoders bridge two questions: SAEs
ask ``what features are present in an activation?''; transcoders ask
``what features does this MLP \emph{compute}, and how do they depend
on its inputs?''. The latter is the question circuit analysis needs
answered.

\subsection{Forward link: are these features canonical?}
\label{sec:sae-forward}

The architectures of this section --- ReLU+$\ell_1$, Top-K, JumpReLU,
and transcoders --- give us a working pipeline for converting
polysemantic neuron-basis activations into a sparse, named feature
basis at scale. What they do not give us, on their own, is a guarantee
that the basis is unique. Two independent SAE training runs at
different widths $M$ on the same activations may produce different
feature inventories \citep{leask2025-sae-not-canonical}; a feature in
a wider SAE may decompose into a combination of features in a narrower
one (\emph{non-atomicity}); features present in a wider SAE may have no
counterparts in the narrower (\emph{incompleteness}). The
``true features'' framing implicit in
\citet{templeton2024-scaling-monosemanticity} treats SAE features as
the model's underlying units of computation; the
\citet{leask2025-sae-not-canonical} critique argues this is unwarranted.
Section~\ref{sec:sae-canonical} takes up the canonicality question
directly, both for SAEs and for the transcoders that inherit it.
